% Part: lambda-calculus
% Chapter: church-rosser
% Section: parallel-beta-eta-reduction

\documentclass[../../../include/open-logic-section]{subfiles}

\begin{document}

\olfileid[pt]{lam}{cr}{pbe}

\olsection{$\beta\eta$-redução paralela}

Nesta seção, provamos a propriedade de Church-Rosser para a
$\beta\eta$-redução paralela, a noção de redução paralela correspondente à
$\beta\eta$-redução.

\begin{defn}[$\beta\eta$-redução paralela, $\beredpar$] \ollabel{defn:beredpar}
  A \emph{$\beta\eta$-redução paralela} ($\beredpar$) em termos é definida
  indutivamente da seguinte maneira:
  \begin{enumerate}
    \item \ollabel{defn:beredpar1} $x \beredpar x$.
    \item \ollabel{defn:beredpar2} Se $N \xrightarrow{\beta} N'$, então
      $\lambd[x][N] \beredpar \lambd[x][N']$.
    \item \ollabel{defn:beredpar3} Se $P \beredpar P'$ e $Q \beredpar Q'$, então
      $PQ \beredpar P'Q'$.
    \item \ollabel{defn:beredpar4} Se $N \beredpar N'$ e $Q \beredpar Q'$, então
      $(\lambd[x][N])Q \beredpar \Subst{N'}{Q'}{x}$.
    \item \ollabel{defn:beredpar5} Se $N \beredpar N'$, então
      $\lambd[x][Nx] \beredpar N'$, desde que $x \notin FV(N)$.
  \end{enumerate}
\end{defn}

\begin{thm}\ollabel{thm:refl}
  $M \beredpar M$.
\end{thm}

\begin{proof}
  Exercício.
\end{proof}

\begin{prob}
Demonstre \olref[lam][cr][pbe]{thm:refl}.
\end{prob}

\begin{defn}[desenvolvimento $\beta\eta$-completo]\ollabel{defn:becd}
  O \emph{desenvolvimento $\beta\eta$-completo}~$\becd{M}$ de~$M$ é
  definido da seguinte maneira:
  \begin{align}
    \becd{x} &= x \ollabel{defn:becd1} \\
    \becd{(\lambd[x][N])} &= \lambd[x][\becd{N}] \ollabel{defn:becd2}\\
    \becd{(PQ)} &= \becd{P}\becd{Q} && \text{se $P$ não é uma abstração-$\lambd$}
    \ollabel{defn:becd3} \\
    \becd{((\lambd[x][N])Q)} &= \Subst{\becd{N}}{\becd{Q}}{x}
                             \ollabel{defn:becd4} \\
    \becd{(\lambd[x][Nx])} &= \becd{N} \ollabel{defn:becd5} & \text{se $x
                                                      \notin FV(N)$}
  \end{align}
\end{defn}

\begin{lem}\ollabel{lem:comp}
  Se $M \beredpar M'$ e $R \beredpar R'$, então $\Subst{M}{R}{y}
  \beredpar \Subst{M'}{R'}{y}$.
\end{lem}

\begin{proof}
  Por indução na !!{derivation} de $M \beredpar M'$.

  Os quatro primeiros casos são exatamente como os de
  \olref[pb]{lem:comp}. Se a última regra é \olref{defn:beredpar5}, então
  $M$ é $\lambd[x][Nx]$, $M'$ é $N'$ para alguns $x$ e~$N'$, onde
  $x \notin FV(N)$ e $N \beredpar N'$. Queremos mostrar que
  $\Subst{(\lambd[x][Nx])}{R}{y} \beredpar \Subst{N'}{R'}{y}$, isto é,
  $\lambd[x][\Subst{N}{R}{y} x] \beredpar \Subst{N'}{R'}{y}$. Isso decorre
  de \olref{defn:beredpar}\olref{defn:beredpar5} e da hipótese de indução.
\end{proof}

\begin{lem}\ollabel{lem:cont}
  Se $M \beredpar M'$, então $M' \beredpar \becd{M}$.
\end{lem}

\begin{proof}
  Por indução na !!{derivation} de $M \beredpar M'$.

  Os quatro primeiros casos são como os de \olref[pb]{lem:cont}. Se a
  última regra é \olref{defn:beredpar5}, então $M$ é $\lambd[x][Nx]$ e $M'$
  é $N'$ para alguns $x$, $N$, $N'$, onde $x \notin FV(N)$ e
  $N \beredpar N'$. Queremos mostrar que $N' \beredpar
  \becd{(\lambd[x][Nx])}$, isto é, $N' \beredpar \becd{N}$, o que é
  imediato pela hipótese de indução.
\end{proof}

\begin{thm}\ollabel{thm:cr}
  $\beredpar$ tem a propriedade de Church-Rosser.
\end{thm}

\begin{proof}
  Imediato de \olref{lem:cont}.
\end{proof}

\end{document}
